% Guía del instructor de los dos laboratorios SDR, edición 3.0.
% Comparte preámbulo, figuras y banco de preguntas con el manual y los anexos;
% activa las respuestas. Compile los cuatro documentos juntos (compile-xelatex.bat o make).

\documentclass[11pt,a4paper]{scrartcl}
\usepackage{xr-hyper}
\usepackage[palette=signal,mode=manual,body=serif]{marco-academic}
% Preámbulo común del manual del estudiante (main.tex) y de la guía del instructor
% (guia_instructor.tex). Se carga después de \documentclass y de marco-academic.

\usepackage{longtable}
\usepackage{calc}
\usepackage{float}
\usepackage{bookmark}
\usepackage{amssymb}
\usepackage{multirow}
\usepackage{ifthen}
\usepackage{tikz}
\usetikzlibrary{arrows.meta,positioning,calc,decorations.pathreplacing,decorations.pathmorphing,shapes.geometric,fit,backgrounds,babel}

\setcounter{secnumdepth}{2}
\setcounter{tocdepth}{1}
\RedeclareSectionCommand[tocnumwidth=2.5em]{section}
\setlength{\LTleft}{0pt}
% Con \LTright rígido, el \caption de longtable desborda exactamente \linewidth.
\setlength{\LTright}{0pt plus 1fill}
\setlength{\tabcolsep}{4pt}
\renewcommand{\arraystretch}{1.14}

% ---------------------------------------------------------------------------
% Numeración por laboratorio: 0.x (común), 1.x (Windows), 2.x (Ubuntu), A.x.
% Cada laboratorio reinicia secciones y pasos; el prefijo de paso (W/U) evita
% confusiones al citar «el paso W12» frente a «el paso U12».
% ---------------------------------------------------------------------------
\newcommand{\labnum}{0}
\renewcommand{\thesection}{\labnum.\arabic{section}}
\newcounter{paso}
\newcommand{\pasoprefijo}{P}
\renewcommand{\thepaso}{\pasoprefijo\arabic{paso}}
\newcommand{\IniciarLaboratorio}[2]{% #1 número (0,1,2,A), #2 prefijo de paso
  \renewcommand{\labnum}{#1}\setcounter{section}{0}%
  \renewcommand{\pasoprefijo}{#2}\setcounter{paso}{0}}

% Portadilla de laboratorio: aparece en el índice y en los marcadores del PDF.
\newcommand{\PortadillaLaboratorio}[3]{% #1 kicker, #2 título, #3 entradilla
  \clearpage
  \phantomsection
  \addcontentsline{toc}{part}{#2}%
  \begin{tcolorbox}[enhanced,frame hidden,colback=MAPrimaryDark,coltext=white,arc=2mm,
      left=7mm,right=7mm,top=6mm,bottom=6mm]
    {\sffamily\bfseries\footnotesize\color{white!80!MAPrimarySoft}\MakeUppercase{#1}}\par\vspace{2mm}
    {\sffamily\bfseries\fontsize{21}{25}\selectfont #2\par}
  \end{tcolorbox}
  \vspace{2mm}
  \MAIntro{#3}
}

% ---------------------------------------------------------------------------
% Recuadros adicionales sobre el sistema de marco-academic.
% ---------------------------------------------------------------------------
\newtcolorbox{esperado}{ma/base,colback=MASuccessSoft,coltitle=MASuccess,borderline west={2.5pt}{0pt}{MASuccess},title={QUÉ DEBE OBSERVAR Y POR QUÉ}}
\newtcolorbox{alerta}[1][]{ma/base,colback=MADangerSoft,coltitle=MADanger,borderline west={3pt}{0pt}{MADanger},title={SEGURIDAD\ifstrempty{#1}{}{ · #1}}}
\newtcolorbox{diagnostico}[1][]{ma/base,colback=MASurfaceAlt,coltitle=MAInk,borderline west={2.5pt}{0pt}{MAMuted},title={DIAGNÓSTICO\ifstrempty{#1}{}{ · #1}}}
\newtcolorbox{prediccion}[1][]{ma/base,colback=MAWarningSoft,coltitle=MAWarning,borderline west={2.5pt}{0pt}{MAWarning},title={PREDIGA ANTES DE MEDIR\ifstrempty{#1}{}{ · #1}}}
\newtcolorbox{estado}[1][]{ma/base,colback=MATertiarySoft,coltitle=MATertiary,borderline west={2.5pt}{0pt}{MATertiary},title={ESTADO DE VALIDACIÓN\ifstrempty{#1}{}{ · #1}}}

% Paso detallado: acción, dónde, parámetro, por qué, observación, verificación, si falla.
% Un argumento vacío omite su fila (no todos los pasos necesitan las siete partes).
\newtcolorbox{pasobox}[2]{ma/base,breakable=false,colback=MASurface,colframe=MARule,frame hidden,
  borderline={0.4pt}{0pt}{MARule},borderline west={2.5pt}{0pt}{MAPrimary},
  coltitle=MAPrimaryDark,title={PASO #1\enspace·\enspace#2},left=3.5mm,right=3.5mm,top=2.2mm,bottom=2.4mm,
  before skip=6pt,after skip=6pt}
\newcommand{\MApasofila}[2]{\ifstrempty{#2}{}{\textcolor{MAMuted}{\sffamily\bfseries\scriptsize #1} & #2\\}}
\newcommand{\Paso}[7]{%
  \refstepcounter{paso}%
  \begin{pasobox}{\thepaso}{#1}
    \small
    \begin{tabularx}{\linewidth}{@{}>{\raggedright\arraybackslash}p{2.3cm}>{\raggedright\arraybackslash}X@{}}
      \MApasofila{DÓNDE}{#2}%
      \MApasofila{PARÁMETRO}{#3}%
      \MApasofila{POR QUÉ}{#4}%
      \MApasofila{SE OBSERVA}{#5}%
      \MApasofila{VERIFIQUE}{#6}%
      \MApasofila{SI FALLA}{#7}%
    \end{tabularx}
  \end{pasobox}}

% Orden de terminal: comando, propósito, salida esperada, cómo reconocer el fallo, recuperación.
% El comando se escribe tal cual (\detokenize): no use %, # ni llaves desbalanceadas.
\newtcolorbox{comandobox}[1]{ma/base,breakable=false,colback=MASurface,frame hidden,
  borderline={0.4pt}{0pt}{MARule},borderline west={2.5pt}{0pt}{MASecondary},
  left=3.5mm,right=3.5mm,top=2.2mm,bottom=2.4mm,before skip=6pt,after skip=6pt,
  title={\ttfamily\small\color{MAInk}\hyphenpenalty=10000\exhyphenpenalty=10000\$\ #1},
  coltitle=MAInk,fonttitle={},halign title=flush left}
\newcommand{\Comando}[5]{%
  \begin{comandobox}{\detokenize{#1}}
    \small
    \begin{tabularx}{\linewidth}{@{}>{\raggedright\arraybackslash}p{2.3cm}>{\raggedright\arraybackslash}X@{}}
      \MApasofila{PROPÓSITO}{#2}%
      \MApasofila{SE OBSERVA}{#3}%
      \MApasofila{FALLO}{#4}%
      \MApasofila{RECUPERACIÓN}{#5}%
    \end{tabularx}
  \end{comandobox}}

% Rótulos de interfaz y valores: estilo uniforme en todo el manual.
\newcommand{\gui}[1]{\textsf{\textbf{#1}}}
\newcommand{\menu}[1]{\textsf{#1}}
\newcommand{\valor}[1]{\textbf{#1}}
\newcommand{\tecla}[1]{{\setlength{\fboxsep}{1.2pt}\fbox{\sffamily\scriptsize #1}}}
\newcommand{\cmd}[1]{\texttt{\detokenize{#1}}}

% Etiquetas de decisión de preparación de hardware.
\newcommand{\PASS}{{\sffamily\bfseries\color{MASuccess}PASS}}
\newcommand{\COND}{{\sffamily\bfseries\color{MAWarning}CONDICIONAL}}
\newcommand{\BLOQ}{{\sffamily\bfseries\color{MADanger}BLOQUEADO}}

% Casillas para hojas de registro.
\newcommand{\campo}[1][3.2cm]{\rule[-0.4ex]{#1}{0.4pt}}

% ---------------------------------------------------------------------------
% Cuestionarios: una sola fuente (contenido/cuestionarios.tex). El manual del
% estudiante imprime enunciado y opciones; la guía del instructor añade la
% respuesta y su justificación. Activar con \guiainstructortrue.
% ---------------------------------------------------------------------------
\newif\ifguiainstructor
\newcounter{pregunta}
\newcommand{\Pregunta}[7]{% #1 enunciado, #2-#5 opciones A-D, #6 letra correcta, #7 justificación
  \refstepcounter{pregunta}%
  \Needspace{6\baselineskip}%
  \noindent\textbf{\sffamily\small\color{MAPrimaryDark}\thepregunta.}\enspace #1\par\vspace{1pt}
  \begin{enumerate}[label=\textsf{\textbf{\Alph*}}),leftmargin=2.4em,itemsep=0pt,topsep=1pt]
    \item #2
    \item #3
    \item #4
    \item #5
  \end{enumerate}
  \ifguiainstructor
    \vspace{-1pt}{\small\noindent\colorbox{MASuccessSoft}{\sffamily\bfseries\scriptsize\color{MASuccess}\,RESPUESTA #6\,}\enspace #7\par}
  \fi
  \vspace{5pt}}

\setlist[itemize]{leftmargin=1.55em,itemsep=0.25em,topsep=0.35em}
\setlist[enumerate]{leftmargin=1.8em,itemsep=0.35em,topsep=0.35em}

\definecolor{grisclaro}{HTML}{F3F4F6}
\lstset{
	style=ma-code,
	basicstyle=\ttfamily\small,
	backgroundcolor=\color{grisclaro},
	frame=single,
	rulecolor=\color{gray!35},
	breaklines=true,
	columns=fullflexible,
	keepspaces=true,
	showstringspaces=false,
	numbers=none,
	xleftmargin=0.4em,
	xrightmargin=0.4em,
	aboveskip=0.6em,
	belowskip=0.6em
}

% Estilo común de los diagramas TikZ: alto contraste, texto sans legible.
\tikzset{
  bloque/.style={draw=MAPrimaryDark,line width=0.7pt,rounded corners=2pt,fill=MAPrimarySoft,
    align=center,font=\sffamily\scriptsize,minimum height=9mm,inner sep=3pt},
  bloquealt/.style={bloque,draw=MASecondary,fill=MASecondarySoft},
  bloqueinst/.style={bloque,draw=MATertiary,fill=MATertiarySoft},
  bloqueext/.style={bloque,draw=MAMuted,fill=MASurfaceAlt,dashed},
  flecha/.style={-{Stealth[length=2.4mm]},line width=0.8pt,color=MAInk},
  flechadash/.style={flecha,dashed,color=MAMuted},
  rotulo/.style={font=\sffamily\scriptsize,color=MAMuted,align=center},
}

% Banco único de preguntas. \Pregunta{enunciado}{A}{B}{C}{D}{letra}{justificación}.
% El manual del estudiante no imprime la letra ni la justificación (\ifguiainstructor).
% Cada cuestionario reinicia la numeración.

% ---------------------------------------------------------------------------
\newcommand{\CuestionarioDiagnosticoW}{%
\setcounter{pregunta}{0}%
\Pregunta{¿Cuál es la longitud de onda aproximada de una señal de 100 MHz?}
{30 cm.}{3 m.}{30 m.}{300 m.}
{B}{$\lambda = c/f = 3\times10^8/10^8 = 3$ m. Las demás confunden un factor 10.}

\Pregunta{Un nivel de 0 dBm corresponde a:}
{0 W.}{1 W.}{1 mW.}{1 $\mu$W.}
{C}{dBm es potencia referida a 1 mW: $10\log_{10}(P/1\,\mathrm{mW})$; 0 dBm es 1 mW. 0 W sería $-\infty$ dBm.}

\Pregunta{Si la potencia aumenta 3 dB, la potencia lineal:}
{Se duplica aproximadamente.}{Aumenta en 3 mW.}{Se triplica.}{Aumenta un 3~\%.}
{A}{$10^{0{,}3}\approx 2$. Los decibelios expresan cocientes, no sumas de potencia absoluta.}

\Pregunta{En una señal FM, la información se transporta en:}
{La amplitud de la portadora.}{La polarización de la onda.}{La posición de pulsos en el tiempo.}{La frecuencia instantánea de la portadora.}
{D}{$f_i(t) = f_c + k_f m(t)$. La amplitud es la de AM; la posición de pulsos es la PPM del Laboratorio 02.}

\Pregunta{El ancho de banda de una señal es:}
{El intervalo de frecuencias que ocupa.}{Su frecuencia central.}{Su amplitud máxima.}{Su velocidad de propagación.}
{A}{Una emisora de FM centrada en 100,9 MHz ocupa unos 200 kHz alrededor de ese centro: centro y ancho son conceptos distintos.}

\Pregunta{Con muestreo complejo (I/Q), ¿qué tasa mínima permite observar 2,4 MHz de espectro a la vez?}
{1,2 MSa/s.}{2,4 MSa/s.}{4,8 MSa/s.}{24 MSa/s.}
{B}{Con muestras complejas la ventana observable es igual a la tasa de muestreo (de $-f_s/2$ a $+f_s/2$ alrededor del centro). El criterio $2B$ se aplica al muestreo real.}

\Pregunta{En su modo normal, un analizador de espectro muestra:}
{Fase frente a frecuencia.}{Audio demodulado.}{Amplitud frente a tiempo únicamente.}{Potencia frente a frecuencia.}
{D}{Es su función básica. Solo en span cero muestra amplitud frente a tiempo.}

\Pregunta{Una antena especificada para 868 MHz se conecta a un receptor sintonizado en 100 MHz. Lo correcto es afirmar que:}
{No puede recibir nada en ningún caso.}{Rinde igual que en 868 MHz.}{Puede recibir señales intensas, pero su rendimiento a 100 MHz no está especificado.}{Convierte la señal de 100 MHz a 868 MHz.}
{C}{Fuera de banda la antena sigue captando campo, con eficiencia desconocida y normalmente baja. Una antena no cambia la frecuencia de la señal.}
}

% ---------------------------------------------------------------------------
\newcommand{\CuestionarioFinalW}{%
\setcounter{pregunta}{0}%
\Pregunta{¿Por qué el demodulador de FM de radiodifusión usa unos 180--200 kHz y el de FM estrecha unos 12,5 kHz?}
{Porque la radiodifusión usa más potencia.}{Porque la radiodifusión tiene mayor desviación ($\pm$75 kHz) y una banda base más ancha (MPX hasta unos 57 kHz).}{Porque 100 MHz es una frecuencia más alta que la de las radios portátiles.}{Porque el HackRF exige ese ancho mínimo.}
{B}{Carson: $B\approx 2(\Delta f + f_m)$. La potencia y la frecuencia portadora no determinan el ancho ocupado; el HackRF no impone el ancho del canal.}

\Pregunta{El dispositivo está en 101,2 MHz y la emisora en 100,9 MHz. El desplazamiento del canal debe ser:}
{+300 kHz.}{0 Hz.}{$-300$ kHz.}{$-101{,}2$ MHz.}
{C}{$\Delta f = f_{\mathrm{emisora}} - f_{\mathrm{dispositivo}} = 100{,}9 - 101{,}2 = -0{,}3$ MHz.}

\Pregunta{¿Para qué se sintoniza el dispositivo 300 kHz por encima de la emisora?}
{Para no demodular sobre el pico DC del centro del espectro.}{Para aumentar la potencia recibida.}{Porque la normativa lo exige.}{Para reducir la tasa de muestreo necesaria.}
{A}{El receptor de conversión directa genera un artefacto en el centro; el desplazamiento aleja la emisora sin perderla. No cambia la potencia ni la tasa.}

\Pregunta{Con una tasa de 2,4 MSa/s y decimación 1, SDRangel muestra a la vez:}
{1,2 MHz de espectro.}{4,8 MHz de espectro.}{Toda la banda de 88--108 MHz.}{2,4 MHz de espectro.}
{D}{Muestreo complejo: ventana igual a la tasa. Por eso se recorre la banda en pasos de 2 MHz.}

\Pregunta{Al pasar LNA/VGA de 16/16 a 40/40 dB con emisoras locales intensas, lo más probable es que:}
{La relación señal/ruido mejore siempre.}{Suban señal y ruido y aparezcan espurios por saturación.}{Aumente la potencia que transmite la emisora.}{Se estreche el ancho de banda de la emisora.}
{B}{La ganancia amplifica todo; con señales fuertes el receptor se satura y crea productos de intermodulación. La ganancia del receptor no altera el transmisor.}

\Pregunta{¿Por qué se mantiene \gui{RF amp} apagado al empezar?}
{Porque añade 14 dB delante de todo y facilita la saturación y el riesgo con señales intensas.}{Porque solo sirve para transmitir.}{Porque cambia la frecuencia del dispositivo.}{Porque desactiva el filtro BBF.}
{A}{Es la primera etapa de ganancia: con él activo, una señal fuerte satura antes el receptor. La documentación del HackRF advierte además que activar el amplificador por error con señales intensas puede causar daño permanente. También se usa en recepción.}

\Pregunta{¿Qué aporta la cascada que no aporta el espectro instantáneo?}
{Mayor resolución en frecuencia.}{El nivel absoluto en dBm.}{La evolución temporal de cada frecuencia, con la potencia codificada en color.}{El audio demodulado.}
{C}{La cascada apila espectros en el tiempo: distingue emisiones continuas de intermitentes. No da dBm absolutos ni audio.}

\Pregunta{En el analizador, sin señal en la pantalla, se pasa la RBW de 100 kHz a 10 kHz. ¿Qué ocurre con el piso de ruido mostrado y con el tiempo de barrido automático?}
{El piso baja unos 10 dB y el barrido se alarga.}{El piso sube unos 10 dB y el barrido se acorta.}{El piso no cambia y el span se multiplica por 10.}{El piso baja 1 dB y el barrido no cambia.}
{A}{La potencia de ruido que atraviesa el filtro es proporcional a la RBW: $10\log_{10}(100/10) = 10$ dB menos. Un filtro más estrecho tarda más en responder, por eso el barrido automático se alarga. La RBW no cambia el span.}

\Pregunta{Una VBW menor que la RBW sirve para:}
{Aumentar la resolución en frecuencia.}{Proteger la entrada del analizador.}{Cambiar la frecuencia central.}{Suavizar la traza (promediado de vídeo) sin cambiar la resolución.}
{D}{La VBW filtra después del detector; la resolución la fija la RBW. La protección de entrada la da la atenuación.}

\Pregunta{El SIGLENT marca $-45$ dBm en la emisora y SDRangel muestra el pico a $-20$ dB. Lo correcto es:}
{El HackRF es 25 dB más sensible.}{El analizador está averiado.}{Las cifras no son comparables: SDRangel usa dB relativos al fondo de escala digital y el SIGLENT dBm absolutos en su entrada.}{Basta sumar 25 dB a SDRangel para calibrarlo.}
{C}{Distintas referencias, ganancias y RBW. Una única comparación no calibra nada.}

\Pregunta{¿Qué observación demuestra que se recuperó información (nivel 2) y no solo energía?}
{Un pico en el analizador.}{Una traza continua en la cascada.}{Un aumento de la potencia del canal.}{Audio inteligible a la salida del demodulador.}
{D}{Picos, trazas y potencia son nivel 1: existe energía. Solo la moduladora recuperada demuestra demodulación.}

\Pregunta{En Colombia, el demodulador quedó con de-énfasis de 50 $\mu$s en lugar de 75 $\mu$s. Lo esperable es:}
{Que no haya audio.}{Que la frecuencia de la emisora se desplace.}{Que el audio suene con los agudos realzados, porque la compensación del pre-énfasis es incompleta.}{Que se pierda la señal piloto de 19 kHz.}
{C}{El transmisor realza los agudos con 75 $\mu$s; deshacerlo con 50 $\mu$s los atenúa menos de lo necesario. El audio existe y la portadora no cambia.}
}

% ---------------------------------------------------------------------------
\newcommand{\CuestionarioDiagnosticoU}{%
\setcounter{pregunta}{0}%
\Pregunta{En el radar secundario, las interrogaciones y las respuestas usan:}
{1090 MHz para interrogar y 1030 MHz para responder.}{978 MHz en ambos sentidos.}{1030 MHz para interrogar y 1090 MHz para responder.}{1575 MHz en ambos sentidos.}
{C}{Enlace ascendente 1030 MHz, descendente 1090 MHz. 978 MHz es UAT; 1575 MHz es GPS L1.}

\Pregunta{La diferencia esencial entre radar primario y secundario es que el primario:}
{Detecta el eco reflejado sin necesidad de equipo a bordo.}{Necesita un transpondedor en la aeronave.}{Solo funciona con ADS-B.}{Transmite la identificación de la aeronave.}
{A}{El primario es no cooperativo; el secundario depende del transpondedor e identifica.}

\Pregunta{¿Cuántas direcciones distintas permite una dirección ICAO de 24 bits?}
{24.}{4096.}{Unos 1000 millones.}{Unos 16,8 millones.}
{D}{$2^{24} = 16\,777\,216$. 4096 son los códigos del modo A ($2^{12}$).}

\Pregunta{A 1 Mbit/s, cada bit dura:}
{1 ms.}{1 $\mu$s.}{1 ns.}{0,5 s.}
{B}{$T = 1/R = 10^{-6}$ s.}

\Pregunta{El byte hexadecimal \texttt{8D} en binario es:}
{\texttt{10001101}.}{\texttt{10011000}.}{\texttt{11010001}.}{\texttt{10001011}.}
{A}{8 = \texttt{1000}, D = \texttt{1101}.}

\Pregunta{Una comprobación de paridad o CRC en un mensaje digital sirve para:}
{Cifrar el contenido.}{Detectar (y a veces corregir) errores de transmisión.}{Comprimir el mensaje.}{Aumentar el alcance de la señal.}
{B}{Es redundancia calculada sobre los datos; no oculta ni comprime nada.}

\Pregunta{La multilateración (MLAT) necesita:}
{Un único receptor con buena antena.}{Que la aeronave transmita su posición GNSS.}{Varios receptores sincronizados que midan diferencias de tiempo de llegada.}{Un radar primario.}
{C}{La posición sale de diferencias de tiempo entre receptores; un solo receptor no basta.}

\Pregunta{En Ubuntu, \texttt{sudo} antes de una orden:}
{La ejecuta con privilegios de administrador.}{La ejecuta en segundo plano.}{Instala el programa.}{Muestra la ayuda de la orden.}
{A}{Por eso no se usa para programas gráficos como SDRangel.}
}

% ---------------------------------------------------------------------------
\newcommand{\CuestionarioFinalU}{%
\setcounter{pregunta}{0}%
\Pregunta{¿Por qué observar energía en 1090 MHz no demuestra todavía una trama Mode S válida?}
{Porque la energía puede ser ruido u otros emisores, y solo la paridad coherente con el formato valida los bits.}{Porque el HackRF no recibe en 1090 MHz.}{Porque las tramas válidas no producen energía visible.}{Porque solo el SIGLENT puede validar tramas.}
{A}{Energía es nivel 1; la trama válida es nivel 4. El analizador no decodifica.}

\Pregunta{Los pulsos del preámbulo Mode S comienzan en:}
{0; 0,5; 1,0 y 1,5 $\mu$s.}{0; 2; 4 y 6 $\mu$s.}{0; 1,0; 3,5 y 4,5 $\mu$s.}{0; 8; 16 y 24 $\mu$s.}
{C}{Cuatro pulsos de 0,5 $\mu$s en 0; 1,0; 3,5 y 4,5 $\mu$s; los datos empiezan en 8 $\mu$s.}

\Pregunta{En la PPM de Mode S, un bit «1» se representa con:}
{Un pulso en la segunda mitad del intervalo de 1 $\mu$s.}{Un pulso en la primera mitad del intervalo de 1 $\mu$s.}{Ausencia de pulso.}{Dos pulsos en el intervalo.}
{B}{«1» = pulso en la primera mitad; «0» = en la segunda. Siempre hay un pulso por bit.}

\Pregunta{Una DF17 se valida sola y una DF4 no. ¿Por qué?}
{La DF4 no lleva paridad.}{La DF17 se transmite con más potencia.}{La DF4 usa otro polinomio.}{En la DF17 el resto de la paridad debe ser cero; en la DF4 el resto es la dirección ICAO, que hay que conocer de antemano.}
{D}{Mismo polinomio, distinta superposición: PI frente a AP. Una AP «válida» con una dirección desconocida puede ser ruido.}

\Pregunta{Una trama DF17 con \emph{Type Code} 4 transporta:}
{La identificación (indicativo) de la aeronave.}{La posición en vuelo.}{La velocidad.}{El código de modo A.}
{A}{TC 1--4: identificación; 9--18: posición en vuelo con altitud barométrica; 19: velocidad.}

\Pregunta{Una aeronave aparece en la tabla con altitud, pero no en el mapa. La causa más probable es:}
{El demodulador está mal sintonizado.}{Aún no hay un par par/impar reciente ni una referencia que permita calcular la posición CPR.}{La aeronave no tiene dirección ICAO.}{El mapa solo muestra aeronaves con indicativo.}
{B}{La posición es información derivada; exige mensajes de posición suficientes y una referencia adecuada.}

\Pregunta{La referencia del receptor está en unos 4,7°N y la aeronave real en 37,1°N. Con decodificación CPR local, el resultado será:}
{Correcto, porque la trama incluye la latitud completa.}{Un error de configuración de SDRangel.}{Una posición falsa en la zona de latitud más cercana a la referencia, sin aviso.}{Imposible de calcular: el mapa queda vacío.}
{C}{La trama codifica la posición dentro de una zona de unos 6°; la decodificación local elige la zona más cercana a la referencia. Así ocurrió en la verificación del respaldo (bloque H del Laboratorio 02).}

\Pregunta{En el respaldo offline, ¿por qué el número de tramas por segundo no describe el tráfico real?}
{Porque SDRangel no cuenta tramas en archivos.}{Porque el archivo está a 2 MSa/s.}{Porque las tramas de la grabación no son válidas.}{Porque el archivo comprime el tiempo (silencios eliminados) y se reproduce en bucle.}
{D}{Las tramas son válidas, pero su ritmo en el archivo no es el del aire.}

\Pregunta{¿Por qué el demodulador ADS-B exige al menos 2 MSa/s?}
{Porque los chips de 0,5 $\mu$s necesitan al menos una muestra por medio bit.}{Porque 1090 MHz dividido entre 2 da 545 MHz.}{Porque es la tasa mínima de cualquier HackRF.}{Porque el mapa necesita esa tasa.}
{A}{2 Mchip/s $\Rightarrow$ 2 MSa/s como mínimo; 4 MSa/s da dos muestras por chip.}

\Pregunta{El modelo y la matrícula que muestra la tabla provienen de:}
{Los mensajes DF17 de identificación.}{El cálculo CPR.}{Una base de datos externa descargada (nivel 7).}{La señal de 1030 MHz.}
{C}{La aeronave transmite dirección e indicativo, no su modelo ni su matrícula.}

\Pregunta{Bajar el umbral de correlación de 7 dB a 3 dB produce:}
{Menos tramas y menos carga de CPU.}{Más intentos de decodificación, más tramas inválidas y más carga de CPU.}{Mayor alcance físico de la antena.}{Ningún cambio.}
{B}{Así lo indica la documentación del plugin: más tramas recuperadas a costa de más errores.}

\Pregunta{¿Por qué un barrido normal del analizador rara vez muestra mensajes de 1090 MHz?}
{Porque el analizador no cubre 1090 MHz.}{Porque las ráfagas duran microsegundos y el oscilador local solo pasa por 1090 MHz una pequeña fracción de cada barrido.}{Porque ADS-B se transmite en polarización horizontal.}{Porque la RBW mínima es demasiado ancha.}
{B}{El SSA3032X Plus llega a 3,2 GHz. El problema es temporal: por eso se usan Max Hold o span cero con disparo de vídeo.}
}

\guiainstructortrue
\externaldocument[M-]{build/laboratorio_sdr_surveillance}[laboratorio_sdr_surveillance.pdf]
\externaldocument[A-]{build/anexo_a_fundamentos_sdr}[anexo_a_fundamentos_sdr.pdf]
\externaldocument[B-]{build/anexo_b_evaluacion_sdr}[anexo_b_evaluacion_sdr.pdf]

\MASetup{
 title={Guía del instructor:\\laboratorios SDR},
 subtitle={Preparación del banco, ensayo previo, conducción de la sesión, claves de los cuestionarios y de las hojas de registro del Laboratorio 01 (FM, Windows) y del Laboratorio 02 (ADS-B y Mode S, Ubuntu)},
 author={Marco Teran},
 institution={Documento para el instructor · no distribuir a estudiantes},
 course={Vigilancia aeronáutica · Instrumentación RF y SDR},
 code={LAB-SDR-1090 · GUÍA v3.0},
 version={3.0},
 date={16 de octubre de 2026},
 header-left={Guía del instructor · Laboratorios SDR},
 header-right={Contiene respuestas},
 cover-note={\textbf{Propósito:} que cada sesión empiece con el banco listo y que la evaluación premie la interpretación honesta de la evidencia.\\[1.5mm]
  \textbf{Contiene} las claves de los cuatro cuestionarios y las respuestas de referencia de las hojas del Anexo B: no la publique junto al material del estudiante.}
}

\hypersetup{
	pdftitle={Guía del instructor: laboratorios SDR},
	pdfsubject={Preparación, conducción de la sesión y claves de evaluación},
	pdfkeywords={SDR, HackRF One, SDRangel, ADS-B, Mode S, FM, SIGLENT, guía del instructor},
	pdflang={es},
	bookmarksnumbered=true,
	bookmarksopen=true
}
\AtBeginDocument{\hypersetup{pdftitle={Guía del instructor: laboratorios SDR},
	pdfsubject={Preparación, conducción de la sesión y claves de evaluación}}}

\begin{document}
\MAMakeTitle

\tableofcontents
\clearpage

\IniciarLaboratorio{I}{I}
% Guía del instructor, edición 3.0 (banco completo).

\section{Resumen operativo}

\begin{resumenejecutivo}
\begin{itemize}
  \item \textbf{Dos laboratorios de solo recepción}: Lab.\ 01 (Windows, FM, 210 min) y Lab.\ 02 (Ubuntu 24.04, ADS-B/Mode S, 225 min).
  \item \textbf{Cuatro documentos}: el \emph{manual} (procedimiento), el \emph{Anexo A} (fundamentos y diagnóstico), el \emph{Anexo B} (cuestionarios y hojas de registro, que el estudiante imprime) y esta guía.
  \item \textbf{Banco completo}: HackRF One, ANT500 para FM, antena de 1090 MHz con base magnética, filtro pasivo de 1090 MHz, cable de extensión SMA, adaptadores SMA, adaptador N -- SMA, SIGLENT SSA3032X Plus y atenuador SMA de 20 dB. La lista con enlaces de compra está en el README del repositorio.
  \item \textbf{Antes de cada curso}: instalar el software y hacer el ensayo previo de la sección~\ref{sec:i-ensayo}.
\end{itemize}
\end{resumenejecutivo}

\section{Requisitos del curso y del estudiante}

\begin{itemize}
  \item \textbf{Conceptuales:} frecuencia y longitud de onda, dB y dBm, ancho de banda, muestreo, AM/FM, radar primario y secundario, binario y hexadecimal. Los cuestionarios diagnósticos detectan carencias: si más de un tercio del grupo falla las preguntas 1--3 del diagnóstico del Lab.\ 01, dedique 10 minutos a dB y longitud de onda durante la orientación.
  \item \textbf{Prácticos:} copiar órdenes con exactitud, abrir una terminal y capturar pantalla. No se asume experiencia con SDR ni con Linux.
  \item \textbf{Organización:} grupos de 2 por banco. Con un solo SIGLENT, rote la etapa 5 del Lab.\ 01 con la etapa 6, y la etapa 8 del Lab.\ 02 con la etapa 7.
\end{itemize}

\section{Tiempo del instructor}

\begin{longtable}{@{}>{\raggedright\arraybackslash}p{0.55\linewidth}>{\raggedleft\arraybackslash}p{0.15\linewidth}>{\raggedright\arraybackslash}p{0.22\linewidth}@{}}
\toprule
\textbf{Actividad} & \textbf{Tiempo} & \textbf{Cuándo} \\
\midrule
\endhead
Instalación de software por puesto & 30--45 min & Una vez por curso \\
Ensayo previo del Lab.\ 01 & 45 min & Antes del curso \\
Ensayo previo del Lab.\ 02 (vivo y grabación) & 60 min & Antes del curso, a la hora de la clase \\
Comprobación rápida por puesto & 10 min & El día de la sesión \\
\bottomrule
\end{longtable}

\section{El banco y su preparación}

\begin{checklist}
  \item HackRF One con su cable USB de datos; los tres SMA limpios.
  \item ANT500 extensible sin holguras; antena de 1090 MHz con su cable de 3 m y base magnética.
  \item Filtro de 1090 MHz: pasivo (no requiere alimentación); géneros de sus conectores compatibles con la cadena, con los adaptadores SMA del paquete del HackRF.
  \item Adaptador N macho -- SMA hembra en cada analizador.
  \item Atenuador SMA de 20 dB guardado junto al banco (solo para emplazamientos muy cercanos a radares; véase la sección~\ref{sec:i-siglent}).
  \item SIGLENT: \gui{Preset} correcto, TG apagado y fecha de calibración anotada.
  \item Una ventana con cielo despejado por aula para la antena de 1090 MHz, con un recorrido de cable seguro.
  \item Grabación de referencia copiada en la carpeta personal de cada usuario de Ubuntu.
\end{checklist}

\section{Software en Windows}

\begin{longtable}{@{}>{\raggedright\arraybackslash}p{0.30\linewidth}>{\raggedright\arraybackslash}p{0.62\linewidth}@{}}
\toprule
\textbf{Componente} & \textbf{Instalación y comprobación} \\
\midrule
\endhead
SDRangel 7.27.2 & Instalador oficial \texttt{sdrangel-7.27.2-win64.exe} (\url{https://github.com/f4exb/sdrangel/releases}), SHA-256 \texttt{96e4d6c83ad554a7 8173310f78302162 fd52463d6a890da5 8211021dabf74dc2}. Incluye el runtime de Visual C++. \\
Herramientas HackRF & radioconda (\url{https://github.com/radioconda/radioconda-installer}), recomendado por la documentación de HackRF; aporta \texttt{hackrf\_info}. \\
Controlador & «HackRF One» sin advertencias en el Administrador de dispositivos. Zadig solo si falta (procedimiento en el Anexo A). \\
Comprobación & \tecla{Ctrl}+\tecla{R} en SDRangel muestra el HackRF con su número de serie. \\
\bottomrule
\end{longtable}

\section{Software en Ubuntu 24.04}

\begin{enumerate}
  \item \texttt{sudo apt update \&\& sudo apt install hackrf} (paquete 2023.01.1-9build1, \emph{universe}); confirme la regla \texttt{/usr/lib/udev/rules.d/60-libhackrf0.rules}.
  \item \texttt{hackrf\_info} sin \texttt{sudo} debe encontrar el HackRF; si no, siga el Anexo A (permisos).
  \item \texttt{sudo snap install sdrangel} y \texttt{sudo snap connect sdrangel:raw-usb}. No use el \texttt{.deb} de 24.04 de la versión 7.27.2: depende de \texttt{libuhd4.10.0}, que no está en los repositorios estándar.
  \item Copie \texttt{datos\_adsb/} a la carpeta personal de cada usuario y ábrala una vez con \textbf{FileInput} para confirmar el acceso desde el snap.
\end{enumerate}

\section{El SIGLENT y los emplazamientos cercanos a radares}
\label{sec:i-siglent}

\begin{itemize}
  \item Límites (hoja de datos DS0703P): RF INPUT N hembra de 50 $\Omega$; potencia media máxima de 30 dBm durante 3 min con $f_c \ge 10$ MHz, atenuación $> 20$ dB y preamplificador apagado; daño a 33 dBm; DC máximo $\pm 50$ V. Con antenas pasivas el riesgo real es la descarga electrostática.
  \item Rótulos de las teclas, \gui{N dB BW}, \gui{Demod} y tiempo de barrido mínimo (900 $\mu$s): manual de usuario UM0703P-E01A. Las funciones \emph{Occupied BW} y \emph{Spectrum Monitor} pertenecen al kit opcional AMK.
  \item \textbf{Radares cercanos.} El HackRF no tiene preselección. A menos de unos pocos kilómetros de un SSR o un DME, el pico recibido puede acercarse a $-5$ dBm (estimación en el Anexo A). En ese caso prohíba el \gui{RF amp}, no oriente la antena hacia el radar y, si el emplazamiento es muy cercano, intercale el atenuador de 20 dB entre el filtro y el HackRF, informando al grupo de que perderá sensibilidad.
\end{itemize}

\section{HackRF: firmware}

Anote con \texttt{hackrf\_info} la versión de firmware y la del software. Herramientas oficiales actuales: v2026.01.3; Ubuntu 24.04 empaqueta 2023.01.1. Si algún programa advierte de incompatibilidad, actualice el firmware siguiendo la documentación oficial de HackRF, fuera de clase: una actualización interrumpida deja la placa en modo DFU.

\section{Versiones de referencia}

\begin{longtable}{@{}>{\raggedright\arraybackslash}p{0.30\linewidth}>{\raggedright\arraybackslash}p{0.62\linewidth}@{}}
\toprule
\textbf{Componente} & \textbf{Versión} \\
\midrule
\endhead
SDRangel & 7.27.2 (Windows: instalador oficial; Ubuntu: snap) \\
Herramientas HackRF & 2026.01.3 (radioconda) / 2023.01.1 (Ubuntu) \\
Sistemas & Windows 10/11; Ubuntu 24.04 LTS nativo \\
Verificador de la grabación & Python 3.12 o superior con \texttt{numpy} \\
\bottomrule
\end{longtable}

\section{Ensayo previo}
\label{sec:i-ensayo}

Siga el manual \emph{tal como está escrito}, en el aula y a la hora de la clase. Anote los resultados en la tabla y corrija el manual si su versión de software muestra algo distinto.

\begin{longtable}{@{}>{\raggedright\arraybackslash}p{0.50\linewidth}>{\raggedright\arraybackslash}p{0.42\linewidth}@{}}
\caption{Registro del ensayo previo.}\label{tab:i-ensayo}\\
\toprule
\textbf{Comprobación} & \textbf{Resultado y fecha} \\
\midrule
\endfirsthead
\toprule
\textbf{Comprobación} & \textbf{Resultado y fecha} \\
\midrule
\endhead
Lab.\ 01, etapas 1--4: audio inteligible de al menos 3 emisoras & \\[2mm]
Lab.\ 01, etapa 5: emisora medida con marcador y $-26$ dB & \\[2mm]
Lab.\ 02, etapas 1--3: tramas válidas en 10 min & \\[2mm]
Lab.\ 02, etapa 6: aeronaves con posición en el mapa & \\[2mm]
Lab.\ 02, etapa 8: tiempo de barrido mínimo en span cero; ráfagas capturadas & \\[2mm]
Lab.\ 02, etapa 9: grabación con valores de referencia & \\[2mm]
Distancia del aula al SSR o DME más cercano & \\[2mm]
\bottomrule
\end{longtable}

\section{Criterios de éxito}

\begin{longtable}{@{}>{\raggedright\arraybackslash}p{0.12\linewidth}>{\raggedright\arraybackslash}p{0.41\linewidth}>{\raggedright\arraybackslash}p{0.39\linewidth}@{}}
\toprule
\textbf{Lab.} & \textbf{Mínimo para aprobar} & \textbf{Excelencia} \\
\midrule
\endhead
01 & HackRF identificado; receptor según la etapa 3; emisora localizada; demodulador con desplazamiento y de-énfasis correctos; audio; analizador con marcador; hoja W8 razonada. & Cuatro experimentos con predicción y explicación física; ancho a $-26$ dB coherente con la regla de Carson; límites de la comparación de niveles bien argumentados. \\
02 & Ubuntu y HackRF verificados; tramas válidas; trama decodificada a mano; tabla con el origen de cada campo; mapa; grabación de referencia con la lección CPR explicada. & Experimento A--B--A concluyente o declarado no concluyente con razones; span cero con duración medida; comparación vivo/grabación. \\
\bottomrule
\end{longtable}

\begin{importante}[Lo que nunca debe premiarse]
Tablas completas con datos que el grupo no pudo observar, capturas de servicios web presentadas como propias o niveles en dBm comparados sin matices con dB de SDRangel.
\end{importante}

\section{Conducción de la sesión}

\begin{longtable}{@{}>{\raggedright\arraybackslash}p{0.22\linewidth}>{\raggedright\arraybackslash}p{0.70\linewidth}@{}}
\toprule
\textbf{Momento} & \textbf{Qué vigilar} \\
\midrule
\endhead
Lab.\ 01, etapa 1 & Antena en CLKIN o CLKOUT: el receptor funciona, pero no se ve ninguna emisora. \\
Lab.\ 01, etapa 3 & Grupos que no cierran radioconda Prompt antes de SDRangel. \\
Lab.\ 01, etapa 4 & Desplazamiento con signo positivo; de-énfasis en 50 $\mu$s; confusión del pico DC con una emisora. Pregunte: «¿qué demuestra ese pico?». \\
Lab.\ 01, etapa 6 & Grupos que cambian dos cosas a la vez; insista en una variable por prueba. \\
Lab.\ 02, etapa 3 & Si a los 10 minutos no hay tramas en ningún puesto, el cielo puede estar vacío: adelante la etapa 9 y vuelva después. \\
Lab.\ 02, etapa 6 & Capturas con coordenadas exactas: recuerde la regla de privacidad. \\
Lab.\ 02, etapa 9 & No anticipe la respuesta: deje que el grupo vea la aeronave en un lugar imposible antes de fijar la referencia. Es el momento de más aprendizaje del laboratorio. \\
\bottomrule
\end{longtable}

\section{Problemas frecuentes}

\begin{longtable}{@{}>{\raggedright\arraybackslash}p{0.30\linewidth}>{\raggedright\arraybackslash}p{0.62\linewidth}@{}}
\toprule
\textbf{Problema} & \textbf{Recuperación} \\
\midrule
\endhead
Ninguna trama en 10 min & Revisar la cadena (antena, filtro, adaptadores); antena en la ventana; LNA 24--32. Si todo está bien, adelantar la etapa 9. No bajar el umbral de 5 dB para «conseguir» filas. \\
Aeronave de la grabación en un lugar imposible & Esperado con la posición del aula. Desactivar \gui{Auto-update position} y fijar 37,0 / 14,0. \\
FileInput no abre el archivo (Ubuntu) & Archivo fuera de la carpeta personal o alterado: copiarlo a \texttt{\textasciitilde/} y comprobar el SHA-256. \\
Hay emisora pero no audio & Signo del desplazamiento, mute, squelch o dispositivo de audio. \\
HackRF «ocupado» & Otro programa abierto. \\
Saturación con RF amp & Apagarlo; aproveche para mostrar los espurios. \\
SDRangel lento & Mapa 2D; no activar \emph{Display All Flight Paths}. \\
\bottomrule
\end{longtable}

\section{La grabación de referencia}

\begin{enumerate}
  \item Integridad: \texttt{sha256sum} debe dar {\ttfamily\small\nolinkurl{d201250f486e5ff14133bc57664c8e5be6c42be3bb989529f29789674ee5f25d}}.
  \item Opcional: \texttt{python3 herramientas/adsb\_iq.py verificar \ldots} informa 430 preámbulos candidatos, 242 tramas válidas y 1 dirección ICAO.
  \item En SDRangel 7.27.2 (Windows), FileInput con bucle y el ADS-B Demodulator con \gui{SR} 2M produjeron la fila \texttt{4d2023}/AMC421: 20\,025--24\,275 ft, 371--390 kt y posiciones 36,96--37,10°N, 13,79--13,86°E con la referencia 37,0/14,0. Con la posición tomada de Windows (Bogotá), la aeronave apareció en Colombia.
  \item La casilla \gui{Auto-update position} viene activada por defecto y sobrescribe cada segundo la posición fijada: desactívela.
\end{enumerate}

\section{Privacidad y uso de datos}

\begin{itemize}
  \item Las capturas publicadas no deben mostrar las coordenadas del receptor.
  \item \gui{Feed} se mantiene desactivado salvo autorización institucional.
  \item La grabación contiene datos públicos radiodifundidos en 2013 por una aeronave real; se usa con fines educativos bajo la licencia BSD-3 de dump1090.
\end{itemize}

\section{Claves de los cuestionarios}
\label{sec:i-claves}

\subsection{Laboratorio 01 · diagnóstico}
\CuestionarioDiagnosticoW
\subsection{Laboratorio 01 · final}
\CuestionarioFinalW
\subsection{Laboratorio 02 · diagnóstico}
\CuestionarioDiagnosticoU
\subsection{Laboratorio 02 · final}
\CuestionarioFinalU

\section{Respuestas de referencia de las hojas}

\begin{longtable}{@{}>{\raggedright\arraybackslash}p{0.17\linewidth}>{\raggedright\arraybackslash}p{0.75\linewidth}@{}}
\toprule
\textbf{Hoja} & \textbf{Respuesta esperada} \\
\midrule
\endhead
W2 & Varias emisoras de alturas distintas: potencia, distancia, obstáculos y respuesta de la antena difieren. \\
W3 & Nivel 2 (audio inteligible); nivel 5 solo si se decodificó RDS. \\
W4 & Ancho a $-26$ dB del orden de 150--300 kHz; nivel en dBm dependiente de la antena y del adaptador. \\
W5 & 50 kHz: distorsión (parte de la desviación fuera del filtro); 100 kHz: muy degradado; 200 kHz: emisora vecina o ruido. \\
W6 & RF BW 80 kHz: distorsión por recorte; 250 kHz: más ruido y posible interferencia adyacente. AF BW 3 kHz: audio «telefónico» pero inteligible. RF BW actúa antes de demodular y AF BW después. \\
W7 & Baja: emisora apenas sobre el ruido. Moderada: mejor compromiso. Alta y excesiva: piso más alto, espurios y distorsión. \\
W8 & Espectro: ambos sí. dBm con exactitud especificada: SIGLENT sí; SDRangel no. Audio: SDRangel sí; SIGLENT también con \gui{Demod} (acepte «sí» o «parcial» si se justifica). Cascada: SDRangel sí; SIGLENT parcial (\emph{Spectrum Monitor} del kit AMK). Ancho con marcadores: SIGLENT sí; SDRangel parcial. 20 MHz: el SIGLENT los recorre en cada barrido, no simultáneamente; el HackRF los vería a la vez a 20 MSa/s, no con 2,4 MSa/s. Antena y cable: ambos sí. \\
W9 & Niveles 1 y 2 (y 5 con RDS); Carson con las medidas de W4; referencias distintas (escala digital frente a dBm), ganancias y RBW distintas, sin calibración. \\
U4 & Un pico es nivel 1: puede ser ruido u otro emisor; solo la paridad coherente con el formato valida los bits. \\
U5 & \texttt{8D} = \texttt{10001101}: DF 17, CA 5. ICAO \texttt{4D2023}. ME \texttt{20} = \texttt{00100000}: TC 4, categoría 0. Grupos: \texttt{000001}, \texttt{001101}, \texttt{000011}, \texttt{110100}, \texttt{110010}, \texttt{110001}, \texttt{100000}, \texttt{100000} = 1, 13, 3, 52, 50, 49, 32, 32 = «AMC421». Es un mensaje de identificación: no contiene posición, y una sola trama de posición tampoco bastaría sin un par o una referencia. \\
U6 & ICAO, indicativo, altitud y velocidad: T. Lat/Lon, distancia y acimut: D. Modelo y matrícula: E. La identificación y la altitud llegan en mensajes distintos, con ritmos distintos. \\
U7 & La posición (decodificación local), la distancia y el alcance dependen de \gui{My Position}; los mosaicos del mapa, de Internet. \\
U8 & Se valora el razonamiento: si A' difiere de A más que B de A, el efecto de la ganancia no está demostrado. El umbral más bajo (C) produce más tramas y más inválidas; sin filtro (D) puede bajar la calidad cerca de emisoras intensas. \\
U9 & El analizador solo ve energía en función del tiempo (nivel 1); no demodula los bits ni comprueba la paridad. \\
U10 & La trama transmite la posición dentro de una zona de unos 6°; la decodificación local eligió la zona más cercana a la referencia del aula. Nada en la trama permite detectar el error. \\
U11 & Niveles 4--6 con evidencia propia; mejoras justificadas con U8 (ubicación antes que ganancia); modelo, matrícula y mosaicos no proceden de la radio. \\
\bottomrule
\end{longtable}

\section{Rúbrica}

La rúbrica para el estudiante está en el Anexo B (sección «Cómo se evalúa»). Para cada criterio, asigne el máximo cuando la evidencia esté registrada, justificada y etiquetada con su nivel; la mitad cuando esté registrada sin justificar; cero cuando falte o no pueda haberse observado.

\section{Limitaciones científicas}

\begin{itemize}
  \item Ningún nivel medido es trazable: antenas y cables no están calibrados; el HackRF no es un instrumento de medida.
  \item La recepción a 1090 MHz depende del lugar, la hora y el tráfico: dos sesiones con el mismo equipo pueden dar resultados muy distintos.
  \item La grabación de referencia tiene el tiempo comprimido y se repite en bucle: no sirve para estudiar tasas, alcance ni seguimiento.
  \item El span cero muestra envolventes, no bits.
  \item El verificador del repositorio es deliberadamente simple: encuentra menos tramas que SDRangel (126 comunes; SDRangel registró 17 DF17 adicionales; el verificador, 1 DF21 que SDRangel no registró).
\end{itemize}


\end{document}
